% =============================================================================
% PQ-Shield journal: drop-in text for the revision.
% Citation keys refer to paper/references.bib. \NUM{...} marks a number to copy
% from paper/key_numbers.md after the clean re-run -- never type numbers by hand.
% =============================================================================
\providecommand{\NUM}[1]{\textbf{[#1]}}

% -----------------------------------------------------------------------------
% §I  Contributions -- replace the streaming bullet so it is positioned against
%     prior art instead of reading as a new construction.
% -----------------------------------------------------------------------------
Fourth, we adapt amortized stream signing~\cite{gennaro1997streams,wong1999flows,perrig2000tesla}
to token-by-token LLM responses, where a post-quantum signature's size (3,309\,B for ML-DSA-65
versus $\approx$71\,B for ECDSA P-256) makes per-chunk signing roughly forty times costlier in bytes.
We compare buffer-and-sign, per-chunk, and hash-chained signing on a real
Llama-3.2-3B-Instruct backend, show that a naive per-chunk design admits an undetected
reordering attack, and close it by binding the chunk index into the signed bytes.

Fifth, we treat per-chunk and hash-chained signing as the two endpoints of one design space:
a hash chain checkpointed every $k$ chunks. We define \emph{exposure} -- the number of tokens a
client displays before any signature covers them -- and measure signature bytes against exposure
and against mid-stream detection of drop and reorder attacks for $k\in\{1,2,5,10,20,\infty\}$,
with an analytic model ($\lfloor n/k\rfloor+1$ signatures, exposure $\le k$ chunks) that the
measurements follow. This gives operators a direct rule for choosing $k$ under an integrity
budget, which matters most for post-quantum signatures: at $k=1$ ML-DSA-65 costs 135,669\,B per
200-token response; at $k=\infty$ only 3,309\,B, but the whole response stays unverified until it ends.

Sixth, we quantify how session resumption amortizes post-quantum key establishment, reporting
end-to-end overhead versus an unprotected control at 1, 10, and 100 requests per key exchange.

% -----------------------------------------------------------------------------
% §II-B  Related work -- three new paragraphs (place after Table I).
% -----------------------------------------------------------------------------
\paragraph*{Post-quantum TLS under network conditions}
Paquin \emph{et al.}~\cite{paquin2020benchmarking} benchmarked post-quantum key exchange and
signatures inside TLS 1.3 using emulated network links (Linux \texttt{netem}), showing that on
lossy or low-bandwidth paths message size, not computation, dominates handshake latency.
Sikeridis \emph{et al.}~\cite{sikeridis2020authentication} reached a similar conclusion for
post-quantum authentication. Montenegro \emph{et al.}~\cite{montenegro2026framework} and
G\'omez-Cambronero \emph{et al.}~\cite{gomezcambronero2026layered} extend this to layered TLS and
application-level load. These works measure the TLS handshake with generic payloads. PQ-Shield
instead measures an application-layer protocol around AI inference payloads, including streaming
responses, and adds controlled adversarial experiments. Our network-emulation experiment
(Section~\ref{sec:netem}) is deliberately coarser than theirs: it shapes the byte stream rather
than packets. We use it to bound, not replace, their packet-level results.

\paragraph*{Stream authentication}
Signing a stream that is not available in full when transmission starts is a classical problem.
Gennaro and Rohatgi~\cite{gennaro1997streams} chain per-block hashes so that a single signature
authenticates an entire stream. Wong and Lam~\cite{wong1999flows} amortize one signature over a
flow with hash trees. TESLA and EMSS~\cite{perrig2000tesla} trade immediate verification for low
per-packet overhead. Our hash-chain strategy is an instance of this family, not a new construction,
and it inherits the same trade-off: integrity is confirmed only when the terminal signature arrives.
What is new here is the setting. LLM token streams have a small, latency-critical first chunk and,
with post-quantum signatures, a per-signature cost large enough that the choice of amortization
strategy decides whether streaming is affordable at all.

\paragraph*{Side channels in encrypted LLM streams}
Weiss \emph{et al.}~\cite{weiss2024keylogging} recovered the content of encrypted AI-assistant
responses from the sizes of streamed packets (the token-length side channel). Their attack is
orthogonal to the choice of key exchange: it applies equally to all four protected configurations
here. It does sharpen one of our qualifications. Per-chunk and hash-chained signing both emit one
record per chunk and therefore expose chunk count and timing, which only buffer-and-sign hides.

% -----------------------------------------------------------------------------
% §IV  System design -- configuration table (replace the three-configuration
%      description) and the naming note.
% -----------------------------------------------------------------------------
\begin{table}[t]
\centering
\caption{Configurations. All protected configurations share AES-256-GCM with an HKDF-SHA256 session key.}
\label{tab:configs}
\setlength{\tabcolsep}{3pt}
\begin{tabular}{@{}llll@{}}
\toprule
ID & Key establishment & Signature & PQ KEX \\
\midrule
Control & none & none & -- \\
A (Classical-RSA) & RSA-2048 OAEP, fresh key per handshake & ECDSA P-256 & no \\
A$'$ (Classical-ECDHE) & ephemeral X25519 & ECDSA P-256 & no \\
B (Hybrid) & ML-KEM-768 & ECDSA P-256 & yes \\
B$'$ (Hybrid-KEX) & X25519MLKEM768 & ECDSA P-256 & yes \\
C (Full PQC) & ML-KEM-768 & ML-DSA-65 & yes \\
\bottomrule
\end{tabular}
\end{table}

Configuration A$'$ is the production classical baseline. TLS 1.3 removed RSA key
transport~\cite{rfc8446}, and ephemeral X25519 is the classical half of the X25519MLKEM768 group
now negotiated by major browsers and CDNs~\cite{ietf-ecdhe-mlkem}. A$'$ and B have identical
signature code, so the A$'\!\rightarrow$B comparison isolates exactly one change: X25519 to
ML-KEM-768. Configuration A generates a fresh RSA-2048 key pair per handshake. No production
system does this, so we report A as a legacy worst case that exposes the cost of RSA key generation
under load, not as \emph{the} classical baseline.

Configuration B pairs a post-quantum key exchange (ML-KEM-768 alone) with a classical signature.
Configuration B$'$ instead combines X25519 and ML-KEM-768 \emph{within} the key exchange, following
the X25519MLKEM768 construction of~\cite{ietf-ecdhe-mlkem}. The server's share is the ML-KEM
encapsulation key followed by an X25519 public key (1,216\,B). The client's share is the ML-KEM
ciphertext followed by its X25519 public key (1,120\,B). The session key is HKDF-SHA256 over both
shared secrets concatenated, so it stays secret unless \emph{both} algorithms are broken. Signature
code is identical, so A$'\!\rightarrow$B$'$ is the migration step deployed on the web today, and
B$\rightarrow$B$'$ is the price of hedging against an ML-KEM break.

% -----------------------------------------------------------------------------
% §IV/§V  Workload split -- one sentence, wherever workloads are introduced.
% -----------------------------------------------------------------------------
Streaming experiments use a real Llama-3.2-3B-Instruct model (Q4\_K\_M, llama.cpp with Metal).
The concurrency, payload-shape, network, HNDL, and MITM experiments use the non-streaming endpoint
with the RandomForest classifier (or the selected payload profile). One locally hosted LLM produces
$\approx$40 tokens per second for a single stream, so a 1,000-connection LLM sweep would
measure the model's queue rather than the cryptography.

% -----------------------------------------------------------------------------
% §V-H  NEW methodology subsection: network emulation.
% -----------------------------------------------------------------------------
\subsection{Network Emulation}\label{sec:netem}
All other experiments run over loopback, where moving an extra 1,088-byte ML-KEM ciphertext or a
3,309-byte ML-DSA signature costs almost nothing. To expose the cost of message size, we place a
TCP proxy between load generator and server that adds a one-way propagation delay of RTT/2 in each
direction and a per-direction bottleneck bandwidth shared by all connections. Bytes are released
in order, so HTTP semantics are unchanged. We evaluate four paths: loopback, \emph{metro}
(10\,ms RTT, 100\,Mbit/s), \emph{WAN} (50\,ms, 20\,Mbit/s), and \emph{mobile} (100\,ms, 5\,Mbit/s),
each with all five configurations at 10 concurrent connections and five repetitions. The proxy
shapes the byte stream above TCP. It does not model slow start, the initial congestion window,
loss, or fragmentation, all of which penalize large handshakes further~\cite{paquin2020benchmarking}.
The measured PQC network cost is therefore a lower bound.

% -----------------------------------------------------------------------------
% §V-I  NEW methodology subsection: checkpointed hash chain and exposure.
% -----------------------------------------------------------------------------
\subsection{Checkpointed Hash Chains and Exposure}\label{sec:checkpoint}
The hash-chain strategy can additionally sign the running chain hash $h_i$ every $k$ chunks.
The client verifies each checkpoint signature against its \emph{own} recomputed $h_i$, so any
dropped, reordered, or altered chunk before a checkpoint makes that checkpoint fail mid-stream.
For an $n$-chunk response this costs $\lfloor n/k\rfloor + 1$ signatures (checkpoints plus the
terminal one), and a decrypted chunk waits at most $k$ chunks for coverage. We record two exposure
measures per transaction: the largest number of chunks displayed but not yet covered by a verified
signature (\emph{max unverified chunks}), and the longest time any chunk waited for coverage
(\emph{verification lag}). Per-chunk signing is the $k=1$ endpoint (exposure one chunk) and plain
hash-chain the $k=\infty$ endpoint (exposure the whole response). We sweep
$k\in\{1,2,5,10,20,\infty\}$ for Hybrid (71\,B signatures) and Full PQC (3,309\,B) on 200-token
Llama responses in 5-token chunks. At each $k$ we repeat the drop and reorder attacks and record
how much of the stream is delivered before detection.

% -----------------------------------------------------------------------------
% §V-J  NEW methodology subsection: session resumption.
% -----------------------------------------------------------------------------
\subsection{Session Resumption}\label{sec:sessions}
By default every transaction performs a fresh key exchange, the worst case. To measure the
amortized case, a client may set \texttt{keep\_session}. The server then caches the derived session
key under the handshake identifier, later requests carry no key-establishment blob and reuse it,
and the last request of the session ends it. Each of 10 concurrent clients reuses one key exchange
for $N\in\{1,10,100\}$ requests (1,000 requests per cell, five repetitions), and we report
end-to-end overhead against Control.

% -----------------------------------------------------------------------------
% §X  Limitations -- replacement list.
% -----------------------------------------------------------------------------
\section{Limitations}
\begin{itemize}
  \item \textbf{Single host.} All measurements come from one machine
  (Apple M3 MacBook Air, 8 cores, 16\,GB, macOS). The load generator and server share its CPU,
  which matters most at 1,000 connections. Relative comparisons between configurations are the
  intended result; absolute latencies are host-specific.
  \item \textbf{Application-layer protocol, not TLS.} PQ-Shield's handshake is a purpose-built
  protocol. Its default fresh key exchange per transaction isolates the asymmetric primitive and is
  a worst case. Section~\ref{sec:sessions} measures the amortized case, but not TLS's own
  resumption mechanisms (PSK, 0-RTT). Server keys are not certificate-authenticated, so it is not a
  substitute for authenticated TLS.
  \item \textbf{Hybrid-KEX measured separately.} Configuration B$'$ was added after the main
  concurrency sweep. Its concurrency and network results come from a supplementary run on the same
  host and day. Its streaming behaviour is not re-measured, because its signature path is identical
  to Configuration B's.
  \item \textbf{Network emulation is byte-level.} The proxy models propagation delay and bandwidth
  but not TCP congestion control, loss, or MTU fragmentation (Section~\ref{sec:netem}). Results are
  a lower bound on post-quantum network cost; packet-level studies~\cite{paquin2020benchmarking,
  sikeridis2020authentication} remain the reference for lossy links.
  \item \textbf{Classical-RSA is a worst case.} Configuration A's per-handshake RSA-2048 key
  generation does not reflect production practice; its availability failure under load should be
  read as a property of that design, contrasted with the realistic Classical-ECDHE baseline.
  \item \textbf{Workloads.} The \texttt{image\_cnn} profile uses an untrained fixed-weight network
  and the \texttt{embedding}/\texttt{llm\_completion} profiles are synthetic. They probe payload
  shape and size, not model behaviour. Streaming uses one quantized 3B-parameter model; the model
  may end a stream before the requested length, so streaming HNDL results are reported against
  tokens actually generated.
  \item \textbf{Timing-validation residual.} A 3--9$\times$ gap between analytically predicted and
  measured streaming \emph{signing time} for some strategy/configuration pairs remains only
  partially explained. Signature-\emph{byte} predictions match exactly and are unaffected.
\end{itemize}
